\chapter{First steps}
\label{ch:first}

\section{Getting nodwyn}

You can download a portable build and both books from the repository's
Gitea Releases page (sign in for this private repository). Choose Linux or
macOS and your computer's architecture, amd64 or arm64. Check the archive's
SHA-256 digest against \file{SHA256SUMS}, extract it, and keep its two
executables beside one another on your \code{PATH}. \code{nodwyn version}
prints the source commit, also recorded in \file{SOURCE.json}. Releases begin
at \code{v0.1.0}; a manual Gitea workflow run on \code{main} supplies the
version and publishes only after checks pass. Pushes dispatch the separate
workstation packages below.

These builds open the browser for \code{nodwyn app}; they do not contain a
native webview window. macOS builds are unsigned and not notarised. Only
Linux amd64 is executed by CI; the other builds are cross-compiled.
\code{nodwyn-egress} is a Linux executable even in the macOS archive,
because it runs inside the MCP container. The release guide at
\file{docs/releases.md} explains downloads and retries. Books are rebuilt
from the same source as the binaries before publication (\adr{42}). Git keeps
the LaTeX sources and catalogue; generated PDFs live in release attachments.
\code{make book} and \code{make expert} also build ignored local copies.

\nodwyn{} is one Go program. From a checkout of the repository:

The private workstation fleet also receives native packages from Beauties.
Each Gitea main push dispatches its exact commit for acceptance, browser checks
and a fake-model selfcheck, then signed Fedora 44 x86-64 RPM and Apple Silicon
Homebrew publication. Routine software updates install the checked build;
\code{nodwyn version} names its source commit. The package includes the native
window and a Linux egress helper. macOS supplies an ad-hoc-signed
\file{Nodwyn.app}. See \file{docs/releases.md}; manual versioned archives,
books and local stable promotion remain separate.

\begin{sh}
make build        # bin/nodwyn, for the terminal
make app          # bin/Nodwyn.app on macOS: the same program with a window
\end{sh}

To build it you need Go; to run it, nothing else: search, the console and
the keychain are all built in. (Two document features call programs you
install yourself: exporting documents needs pandoc, and reading a PDF
attachment needs \code{pdftotext}; Chapter~\ref{ch:workspaces}.) You also
need a \textbf{model} to talk to: \nodwyn{} has no built-in model.
Any Ollama server works, on your machine or another one on your network, and
so does any server that speaks the OpenAI chat-completions protocol, and
Anthropic's own API for Claude.

The native window uses Wails v3, pinned to a beta release
(\adr{43}). The \code{webview} tag and existing console protocol stay the
same. \file{cmd/nodwyn/app\_webview.go} owns the window;
\file{cmd/nodwyn/app\_bridge.go} restricts its three native conveniences
by origin and a per-window capability. Task methods still go through the
ordinary RPC bridge. \code{make app} also creates
\file{bin/Nodwyn-macos-universal.zip}. Native Windows host support and
desktop release builds are not supplied by this shell migration.

\section{The quickest start: the app}

Open \file{Nodwyn.app} (or run \code{nodwyn app}). The first time, Home draws
the forms of Settings \textrightarrow{} Models, one after the other. First the
provider: an Ollama address, Anthropic's API, OpenAI, or any other server from
the presets, and a key for the services that want one. Then the model, from the
list the provider gives, with the efforts it takes; saving it makes it the
default for new tasks. The server is written to a \code{[servers.NAME]}
table of your \file{config.toml}, the model to a \code{[models.NAME]} table
and the choice to \code{[control]}; to choose another model on the same server
later, Settings offers the server by name and the models it lists. A key never
goes into \file{config.toml}. It goes into the encrypted
keychain, which the page makes beside the provider form with a passphrase if there
is none, and a \emph{binding} allows that key only for that server's host.

Then pick a folder, type what you want, and watch. Every task is a chat:
a request that takes long is a message like any other, and the turn goes
on until the model answers or a hard limit stops it. A write in the folder goes ahead, shown in the
conversation with its diff, and the file's previous content is kept so you
can put it back. A command goes to the judge first, if you have one
(Chapter~\ref{ch:keys}). When something needs you (a command the judge did
not allow, or one with no judge to ask, a protected path, content that
looks like a key), the task view shows you the call with the guard's
reasons, and buttons:

\begin{center}
\begin{tabularx}{0.92\linewidth}{cX}
\toprule
\textbf{key} & \textbf{meaning} \\
\midrule
\key{Y} & allow this once \\
\key{T} & allow every write until this turn ends \\
\key{A} & always allow under this folder (or this command prefix), for this task \\
\key{B} & on a keyed request no binding covers: save that binding and allow (only with \code{policy.unbound = "confirm"}) \\
\key{N} & deny: the model is told it was refused and carries on \\
\bottomrule
\end{tabularx}
\end{center}

Not every button appears on every prompt: \key{T} is for writes, which ask
only when you set them to, \key{B} for egress, and some prompts (a write to
\file{.nodwyn/}, content that looks like a key) offer only \key{Y} and
\key{N}. What asks is yours to choose in Settings \textrightarrow{}
Approvals. More on why, and on that setting, in
Chapter~\ref{ch:primitives}.

\section{The terminal way}

Tell \nodwyn{} about your model in \file{config.toml}. \code{nodwyn paths}
prints where it lives:

\begin{txt}
config:   /home/you/.config/nodwyn/config.toml
bindings: /home/you/.config/nodwyn/bindings.toml
keychain: /home/you/.config/nodwyn/keychain
sessions: /home/you/.local/share/nodwyn/sessions
\end{txt}

\begin{toml}
# ~/.config/nodwyn/config.toml
[control]
provider = "ollama"               # or "openai" for any chat-completions server, or "anthropic"
url = "http://localhost:11434"
model = "qwen3:8b"
num_ctx = 81920                   # at least the window plus the reply reserve
\end{toml}

Then go to a folder and start a session:

\begin{txt}
$ cd ~/src/garden
$ nodwyn
control model: http://localhost:11434 qwen3:8b (ollama 0.12.3)
log: /home/you/.local/share/nodwyn/sessions/20260930T200043Z-e0550e21.jsonl
nodwyn session 20260930T200043Z-e0550e21
repo  /home/you/src/garden
tools doc-export fs-edit fs-find fs-list fs-read fs-write index-grep plan-step
      plan-write read-definition second-look shell-exec shell-read shell-send
      shell-start shell-stop
:h for commands

you> what grows in this garden?
\end{txt}

The banner tells you three important things: the \textbf{session id} (the
name of its log, whose path is on the line above), the \textbf{repo} (every
path the model touches is pinned under it), and the \textbf{tools} the model
is offered. You will learn to change that last line in
Chapter~\ref{ch:tools}. (In a git work tree the list also has
\code{git-status} and \code{git-diff}; with keys stored, a few more appear,
Chapter~\ref{ch:keys}.)

\subsection{At the prompt}

\begin{center}
\begin{tabular}{ll}
\toprule
\code{:q} & quit (the log gets a \code{session.end}) \\
\code{:log [n]} & the last \emph{n} log events (default 20) \\
\code{:tools} & the tools offered to the model, with their descriptions \\
\code{:eval EXPR} & evaluate Scheme yourself, with \emph{user} origin \\
\code{:thinking} & show or hide the model's thinking \\
\code{:allow} & the standing approvals in effect \\
\key{Ctrl-C} & abort the running turn; at an empty prompt, exit \\
\bottomrule
\end{tabular}
\end{center}

\subsection{Confirms in the terminal}

The terminal follows the same settings as the console. A write in the
folder goes ahead, and since nobody was asked, its diff is shown dimmed:

\begin{txt}
· ⚑ writes in the folder go ahead (approval.writes)
· --- plants.txt
· +++ plants.txt
· @@ -1,2 +1,3 @@
·  tomatoes
·  basil
· +mint
\end{txt}

The file's previous content is kept first, by git when it has the file
committed and unchanged and by \nodwyn{} otherwise, so it can be put back
(Chapter~\ref{ch:workspaces}). A command no standing approval covers
goes to the judge, if you have one; a call it allows runs, and the log
notes the judged approval. When the judge does not allow it, or there is
no judge, the command stops for you, and the prompt carries a bullet
saying what the judge answered (\emph{``the judge did not allow it:
\ldots''}). With \code{writes = "ask"} in the \code{[approval]} table, a
write stops the same way and shows you the diff:

\begin{txt}
─── confirm: fs-edit ("plants.txt" "basil" "basil\nmint")
  • fs-edit is irreversible
  • changes plants.txt
  --- plants.txt
  +++ plants.txt
  @@ -1,2 +1,3 @@
   tomatoes
   basil
  +mint
allow? [y/N/a=always under plants.txt/t=all writes this turn]
\end{txt}

\code{y} is once; anything else, including just Enter, is no. \code{a}
gives a \textbf{standing approval} for the rest of the session: for a write,
the file's directory (a file at the top of the folder is approved on its
own); for a command, its prefix (\code{go test}, \code{rg}\ldots), which
the judge is then not asked about again. A \code{bash -c} line is offered as
the one command it runs; a line of several commands, or an interpreter
(a shell, \code{python}, \code{env}, \code{ssh}\ldots), offers no \code{a}. \code{t} approves every further
write until the turn ends; each auto-approved diff is still logged and
shown, dimmed.
In the console, \textbf{Don't ask: approve writes and commands} under the
composer goes further for one task: every ordinary write and command is
approved until you uncheck it or the task closes, and the judge is not
asked. Protected paths and writes that look like a key still ask. No Scheme
can turn it on, not even a folder's \file{init.scm}, and a resumed task
starts asking again. The switch is there only once Settings › General › Approvals
offers it (\code{[approval] dont\_ask = true}, off by default, shown with a
red warning); taking the offer away turns it off in every open task.

\begin{tip}
Read the bullet lines. \emph{``command line names a path outside the
repository''} or \emph{``uses the network tool curl''} is the guard telling
you what a yes would do, in plain words. Experts read them every time.
\end{tip}

\section{Try it: your first garden}

\begin{tryit}
Make a folder with a file and ask about it:
\begin{sh}
mkdir -p ~/garden && cd ~/garden
printf 'tomatoes\nbasil\n' > plants.txt
nodwyn
\end{sh}
Then, at the prompt, before asking the model anything:
\begin{txt}
you> :eval (fs-read "plants.txt")
"tomatoes\nbasil\n"
you> :eval (fs-list ".")
("plants.txt")
you> :eval (fs-read "../secret.txt")
⚑ guard deny fs-read: path ../secret.txt is outside the repository ...
error: fs-read: denied: path ../secret.txt is outside the repository ...
\end{txt}
You just called the same primitives the model calls, through the same
guard. The only difference is the origin: yours is \code{user}.
\end{tryit}

\section{Budgets: how a turn stops}

A turn runs until the model answers without asking for a tool, or until it
hits a \textbf{budget}. There are four dimensions, each with a hard limit
(a stop) and a soft one (a checkpoint that asks you whether to continue).
By default each soft limit is its hard one, so no checkpoint interrupts a
turn:

\begin{center}
\begin{tabular}{lll}
\toprule
\textbf{dimension} & \textbf{hard limit} & \textbf{flag for the hard limit} \\
\midrule
tool calls & 500 & \code{-{}-max-tool-calls N} \\
wall time & 4 hours & \code{-{}-budget-wall 2h} \\
model tokens & 10M reported & \code{-{}-budget-tokens N} \\
irreversible actions & 200 attempts & \code{-{}-budget-actions N} \\
spend, the whole task's & none unless you set one & \code{-{}-budget-cost 5} (US dollars) \\
\bottomrule
\end{tabular}
\end{center}

For checkpoints, set a soft limit lower: \code{tool\_calls\_soft},
\code{wall\_soft}, \code{tokens\_soft} or \code{actions\_soft} in
\code{[budget]}, or ``ask at'' in Settings \textrightarrow{} Limits. A soft
limit the file leaves unset follows its hard one, so lowering a hard limit
moves the checkpoint with it. At a checkpoint, \code{y} doubles the soft
limit, \code{n} stops, and a number sets a new soft limit. Approvals and
hard ceilings do not change.

A task started in the console can do without the four turn limits (ADR 49):
untick \textbf{Limits} on Home, or in the task between turns. Its turns then
run, and wait for an answer, as long as they need; the spend limit, the
timeouts and the progress guard still apply. \code{session.info} reports
such a task with \code{limits: false} and a hard limit of 0, none, for the
four, and the log records the choice as \code{session.limits}, which a
resume keeps. A paired device cannot make it.

After each reply the terminal shows what the request used:

\begin{txt}
· [12,400→800 tok · 78% cached · $0.012]
\end{txt}

That is the prompt's tokens, the reply's, the share of the prompt the
provider read from its cache, and what the request cost. The cache share
is the provider's report, or — where it reports none, as older Ollama servers do —
an estimate from the reply's prompt-eval time against the session's
most-cold request, marked ≈; the cost appears only when it is known
(Chapter~\ref{ch:keys}), and
the timings only when the server reports them, as Ollama does. The
request's wall time, as this machine measured it, gives a speed whenever
there are tokens to divide by — \code{22.0~tok/s} — the one speed every
provider gets, prompt processing included. A
checkpoint also says what the turn has cost so far. The token budget adds
up each request's prompt and reply, a prompt token read from the
provider's cache counted as a tenth of one: about what it costs, and what
it asks of the server.

Money has its own limit, if you want one: \code{cost\_soft} and
\code{cost\_hard} in \code{[budget]}, in dollars, or
\code{-{}-budget-cost}. It is the task's, not the turn's: every turn of a
session or a run counts against it, a resumed session's history too, and a
soft limit you raise at its checkpoint stays raised. At the hard limit the
task stops (\code{turn.end} says \code{spend\_limit}), a run ends, and the
next message stops before it reaches the model. Only costs it knows count;
the first turn with a request nobody priced says so in a note.

Independently of budgets, a \textbf{progress guard} ends a turn that goes in
circles: the same call four times in a row, a denied call retried three
times, ten failures in a row, or one tool failing on one file in six of the
last sixteen responses, or a background command that is still running
read thirty times in a row (waiting for one is not repeating a call, so
those reads are counted on their own). In sandbox mode (ADR 55) the same
call four times in a row is not a stop: the fourth result, and every fourth
after it, says so to the model, and the log has a \code{progress warning}
note. No Scheme can catch these stops. That is a
feature: a strategy you write cannot accidentally keep a runaway turn alive.

\section{The scratch directory}

Every session gets a private scratch directory outside your folder, in the
system's temporary directory (so outside \nodwyn's own configuration and
data directories, which the guard protects, and a test suite a command runs
is not refused the temporary files it makes there), for throwaway files: a
script that reproduces a bug, notes, temporary output. The model knows it as \code{\$TMPDIR}: commands run with \code{TMPDIR} set to
it, and a path or command argument starting with \code{\$TMPDIR/} is
rewritten to it before the guard sees it. Writes there need no approval.
Nothing in it reaches your folder or its diff. A clean exit removes it; a
resumed session gets a fresh one.

\begin{tryit}
\begin{txt}
you> :eval (scratch-dir)
"/tmp/nodwyn-scratch-2026...-e0550e21-5431"
\end{txt}
\end{tryit}

\section{Notes that survive the session}

A project can keep notes from one session to the next: how its tests run,
where things are. They live in \nodwyn's data directory, keyed by the
folder's path, \emph{never} in the folder itself:

\begin{scm}
(project-notes)                         ;=> "" the first time
(project-notes-save "Tests: go test ./... ; the fixtures live in testdata/")
\end{scm}

The next session shows them to the model under a heading that calls them
leads to check, not instructions. A save is capped at 8~KiB and redacted,
and anything that looks like a credential needs your yes.

\section{Unattended: \texttt{nodwyn run --yolo}}

A chat already asks little: writes in the folder go ahead, and a judge, if
you have one, answers commands first. But a person is still there for what
the judge does not allow. Sometimes nobody is watching at all: a
benchmark, a CI job, a long task on a throwaway checkout. A run's
\code{-{}-yolo} answers \emph{ordinary} confirms with yes. (In line mode
\code{-{}-yolo} turns on the task's Don't ask, which Settings must offer:
the task is still one a person answers.)
It does not overrule refusals, and it refuses \emph{security} confirms
rather than pretending someone said yes: writes to protected paths, content
that looks like a key, a command line built out of tool output. With nobody
watching, the default loop also changes its manners: it asks the model to
work by evidence and ends a turn only when something changed and a test ran
after it (Chapter~\ref{ch:plans}).

\code{-{}-yolo} still leaves someone to answer a budget checkpoint you set,
and the checkpoints between turns. For a run that nobody answers at
all, and for one whose confirms the judge answers instead of a blanket yes,
see \code{nodwyn run -{}-unattended} and \code{-{}-judge} in
Chapter~\ref{ch:plans}: a batch tool for scripts and benchmarks.

\begin{careful}
The guard makes \code{-{}-yolo} survivable, not safe. Use it on a copy, set
budgets for the job, and read the log afterwards.
\end{careful}

On Linux, \code{python3 scripts/install-desktop.py /absolute/path/to/nodwyn}
installs an application-menu shortcut and the console's icon for your
account. Pin nodwyn from that menu. The native window and launcher share
\code{com.nodwyn.app}; macOS uses the icon in \file{Nodwyn.app}.
